% === Begin Label Data ===
% ID: 406
% 难度星级: 
% 标签: 
% 备注: 
% 组卷引用次数: 0
% === End  Label Data ===

\begin{problem}{2013}{M}{广东卷（文）改}{10}{向量}
设 $ \boldsymbol{a} $ 是已知的平面向量，且 $ \boldsymbol{a}\neq\boldsymbol{0} $．关于向量 $ \boldsymbol{a} $ 的分解，有如下四个命题：  
\circled{1} 给定向量 $ \boldsymbol{b} $，总存在 $ \boldsymbol{c} $，使 $ \boldsymbol{a}=\boldsymbol{b}+\boldsymbol{c} $；  
\circled{2} 给定向量 $ \boldsymbol{b} $ 和 $ \boldsymbol{c} $，总存在实数 $ \lambda $ 和 $ \mu $，使 $ \boldsymbol{a}=\lambda\boldsymbol{b}+\mu\boldsymbol{c} $；  
\circled{3} 给定单位向量 $ \boldsymbol{b} $ 和正数 $ \mu $，总存在单位向量 $ \boldsymbol{c} $ 和实数 $ \lambda $，使 $ \boldsymbol{a}=\lambda\boldsymbol{b}+\mu\boldsymbol{c} $；  
\circled{4} 给定正数 $ \lambda $ 和 $ \mu $，总存在单位向量 $ \boldsymbol{b} $ 和单位向量 $ \boldsymbol{c} $，使 $ \boldsymbol{a}=\lambda\boldsymbol{b}+\mu\boldsymbol{c} $．  
上述命题中的向量 $ \boldsymbol{b} $、$ \boldsymbol{c} $ 和 $ \boldsymbol{a} $ 在同一平面内且两两不共线，则真命题有 \underline{\hspace{4em}}．
\end{problem}

\begin{answer}

\end{answer}

\begin{solutions}

\end{solutions}